\chapter{Design: Joint Application--Network Optimization over Multipath}
\label{ch:p2design}

Part~I chose a single path for a flow. Part~II starts from a more ambitious premise. SCION offers not
one path but several, so why commit to one? And if Mortise~\cite{Shen2026} can tune a
\emph{transport mechanism} to the application, why not tune the \emph{application itself} in the same
breath? This chapter designs a system that (i) uses multiple paths at once via Multipath
QUIC~\cite{DeConinck2017} and (ii) jointly optimizes what the application sends -- the video encoder's
target bitrate -- and how the network spreads it -- the per-path traffic split -- to maximize the
quality of experience of a real-time video call.

\Cref{sec:p2design:motivation} motivates the joint, multipath framing and the strawmen it improves on.
\Cref{sec:p2design:hierarchy} presents the two-timescale hierarchical agent pair.
\Cref{sec:p2design:mdp} specifies the observations, actions, and rewards.
\Cref{sec:p2design:coupling} explains the coupling that keeps the two agents cooperating rather than
working at cross purposes. \Cref{sec:p2design:scoring} extends the scoring idea of Part~I from a discrete
choice to a continuous split over a variable, possibly changing, path set.

\section{Motivation and Strawmen}
\label{sec:p2design:motivation}

Interactive video is a demanding, ubiquitous workload with no playback buffer to hide variance: a frame
that misses its deadline is simply lost. Its QoE depends on two decisions that are usually made in
isolation. The \emph{application} chooses an encoding bitrate -- too high and the network drops frames,
too low and quality suffers. The \emph{network} chooses how to carry the resulting bytes -- over one path
or spread across several. Optimizing either alone is a strawman:

\begin{description}
  \item[Application-only.] Adapt bitrate to the aggregate network but carry it over a fixed or naive
  path assignment. This wastes SCION's path diversity and, as our results show, collapses when the paths
  it depends on degrade.
  \item[Network-only.] Schedule cleverly across paths but treat the offered load as exogenous. The
  scheduler cannot relieve congestion it is itself creating by admitting too high a bitrate.
\end{description}

The two decisions are coupled through the same bottleneck: bitrate sets the load the scheduler must
place, and the scheduler's success sets the effective capacity the bitrate controller sees. Optimizing
them jointly is the design's premise. The challenge is that they naturally operate on \emph{different
timescales} -- an encoder should not change bitrate every frame, but a scheduler must react per frame --
which motivates a hierarchy rather than a single monolithic agent.

\section{Two-Timescale Hierarchical Agents}
\label{sec:p2design:hierarchy}

We split control into two cooperating agents that act at different cadences, a ``brain and body''
decomposition.

\begin{description}
  \item[Application agent (slow).] Every \SI{1}{\second} (once per 30 frames at 30\,fps) it observes
  an aggregate summary of the connection and outputs a single target encoder bitrate in the range
  \SIrange{300}{6000}{\kilo\bit\per\second}. It is deliberately \emph{horizon-agnostic} -- it never
  observes how far along the call is -- so the same policy applies to a call of any length, an essential
  property for a controller meant to run in unbounded real sessions.
  \item[Transport agent (fast).] Every frame (${\approx}\SI{33}{\milli\second}$) it observes per-path
  transport state \emph{and the application agent's current target bitrate}, and outputs a split of that
  frame's bytes across the available paths.
\end{description}

Both are trained with soft actor-critic (\cref{sec:background:rl}), whose continuous action support and
sample efficiency suit these controls and the cost of simulator rollouts.

\section{Observations, Actions, and Rewards}
\label{sec:p2design:mdp}

\paragraph{Application observation and action.} The application agent sees a small fixed-size vector:
normalized current bitrate, smoothed RTT, jitter, loss, and throughput. Its action is one continuous
scalar in $[-1,1]$, mapped affinely to a target bitrate in \SIrange{300}{6000}{\kilo\bit\per\second}.

\paragraph{Transport observation and action.} The transport agent sees a global block (target bitrate,
RTT, loss, throughput) and, for each path, a feature vector (congestion window, smoothed RTT, buffer
occupancy, per-path throughput, per-path loss). Its action is an $N$-vector passed through a softmax to
yield a split $(w_1,\dots,w_N)$ that is non-negative and sums to one -- a valid allocation by
construction.

\paragraph{Rewards.} The rewards encode the QoE objective and the division of responsibility between the
agents. The application's per-window QoE reward composes a quality term with penalties:
\begin{equation}
\mathrm{QoE} = a\cdot\frac{\mathrm{VMAF}(\cdot)}{100}
             - b\cdot\widetilde{\mathrm{lat}}
             - c\cdot\widetilde{\mathrm{jit}}
             - d\cdot\mathrm{loss},
\qquad \text{clipped to } [-2, 1],
\label{eq:p2qoe}
\end{equation}
where $\widetilde{\mathrm{lat}}$ and $\widetilde{\mathrm{jit}}$ are latency and jitter normalized (by
\SI{200}{\milli\second} and \SI{50}{\milli\second}) and clipped, and the default weights are
$(a,b,c,d) = (1, 0.5, 0.5, 1)$. The quality term uses a VMAF surrogate (\cref{ch:p2impl}); when a learned
surrogate that already accounts for loss is used, the explicit loss penalty is dropped to avoid
double-counting. The transport agent receives a per-frame reward that credits delivery without quality,
since it does not control bitrate:
\begin{equation}
R_{\text{transport}} = (1 - \mathrm{loss}) - b\cdot\widetilde{\mathrm{lat}} - c\cdot\widetilde{\mathrm{jit}},
\qquad \text{clipped to } [-2, 1].
\label{eq:p2transport}
\end{equation}

\section{Coupling the Two Agents}
\label{sec:p2design:coupling}

The design's subtlety is keeping the agents cooperative. Two mechanisms couple them.

\paragraph{Observation coupling.} The transport agent observes the application's current target bitrate.
It therefore schedules \emph{knowing} how much load it must place, and can, for instance, spread an
aggressive bitrate across more paths rather than overrunning one. This is what makes the pair hierarchical
rather than two agents optimizing in ignorance of each other.

\paragraph{Reward decomposition.} The application owns quality (the VMAF term) and is credited for the
latency, jitter, and loss experienced over the window its bitrate governed; the transport agent is
credited only for how well it delivered each frame. This assigns each decision the consequences it can
actually control -- the encoder is not blamed for a poor split, nor the scheduler for an over-ambitious
bitrate -- while the shared bottleneck still couples them: a bad split shows up as latency/loss in the
application's reward, and an over-high bitrate shows up as loss the transport agent cannot schedule away.

\section{Continuous Scoring over a Variable Path Set}
\label{sec:p2design:scoring}

The transport agent inherits Part~I's problem -- a variable, possibly changing number of paths -- and its
solution, adapted from a discrete choice to a continuous allocation. A shared per-path encoder maps each
$(\text{global}, \text{path}_i)$ pair to a logit; the logits are softmaxed into the split, so the policy
is permutation-equivariant and independent of the path count. To handle paths that appear and disappear,
each path carries a \emph{liveness} flag and the state includes a binary mask; the softmax excludes dead
paths, so traffic is never allocated to a path that is gone. The critic aggregates per-path embeddings
with a masked mean-pool -- a Deep Sets~\cite{Zaheer2017} operation -- yielding a value estimate invariant
to path ordering and count. This continuous scoring agent is the natural sibling of Part~I's discrete
scoring selector: the same architectural idea, once for ``pick a path'' and once for ``divide across
paths.''
